% =====================================================================
%  CHƯƠNG 5 — HỆ THỐNG TÀI CHÍNH: VN30 VÀ GIÁ VÀNG SJC
% =====================================================================
\chapter{Hệ thống tài chính: VN30 và giá vàng SJC}\label{chap:taichinh}

Phân hệ tài chính trả lời loại câu hỏi mà nhà đầu tư cá nhân Việt Nam hay đặt ra: tỷ giá USD/VND tăng thì cổ phiếu
nào hưởng lợi, Fed tăng lãi suất thì vàng SJC đi đâu, phiên tới FPT có thể ở mức giá nào. Ba câu hỏi đó cần ba loại
tri thức khác nhau. Câu đầu cần \emph{cơ chế} (doanh thu bằng ngoại tệ, chi phí bằng đồng); câu thứ hai cần \emph{suy
luận nhiều bước} qua các biến trung gian (DXY, giá vàng thế giới, giá quy đổi); câu thứ ba cần \emph{dữ liệu chuỗi
thời gian}. Chương này dựng ba lớp tương ứng: \texttt{VietnamFinanceKG} lưu đồ thị nhân quả có dấu và có nguồn,
bộ suy luận đa bước tính chiều tác động ròng trên đồ thị đó, và mạng BiLSTM dự báo lợi suất phiên kế tiếp. Lớp thứ
tư, ít hiển nhiên hơn, là bước \emph{kiểm chứng KG bằng dữ liệu}: dấu của mỗi cạnh lý thuyết được đối chiếu với
tương quan thực nghiệm trên gần 12 năm giao dịch.

\begin{phathien}{Lưu ý về phạm vi}
Hệ thống phục vụ mục đích học tập, không phải công cụ khuyến nghị đầu tư. Kho văn bản tài chính tổng hợp từ văn bản
pháp luật (Nghị định 24/2012/NĐ-CP, Nghị định 232/2025/NĐ-CP), thông cáo của Ngân hàng Nhà nước và báo cáo thường
niên doanh nghiệp; mọi câu trả lời của tác tử đều kết thúc bằng câu ``Thông tin chỉ mang tính tham khảo, không phải
khuyến nghị đầu tư''.
\end{phathien}

% ---------------------------------------------------------------------
\section{Bối cảnh vĩ mô và dữ liệu thị trường}\label{sec:tc-boicanh}
\subsection{Bốn cổ phiếu, bốn kênh truyền dẫn}
Rổ VN30 gồm 30 cổ phiếu vốn hóa và thanh khoản lớn nhất sàn HOSE. Bài chọn bốn mã không phải vì vốn hóa mà vì mỗi
mã đại diện một kênh truyền dẫn vĩ mô khác nhau, để KG có đủ dấu dương và dấu âm cho cùng một biến nguồn:

\begin{itemize}
  \item \textbf{FPT} (công nghệ thông tin): doanh thu xuất khẩu phần mềm thanh toán bằng USD, JPY, EUR, còn chi phí
        chủ yếu là lương kỹ sư trả bằng đồng. Tỷ giá tăng làm doanh thu quy đổi tăng nhanh hơn chi phí, nên về mặt
        cơ bản FPT \emph{hưởng lợi} khi USD/VND tăng; Nhật Bản là thị trường xuất khẩu lớn nhất nên FPT còn nhạy với
        đồng yên.
  \item \textbf{HPG} (thép): nhập quặng sắt và than cốc bằng USD và có nợ vay ngoại tệ, nên tỷ giá tăng làm giá vốn
        tăng và phát sinh lỗ tỷ giá; ngược lại, giải ngân đầu tư công làm tăng cầu thép xây dựng.
  \item \textbf{VCB} (ngân hàng): lợi nhuận phụ thuộc tăng trưởng tín dụng, mà tín dụng nhạy với lãi suất điều hành
        của Ngân hàng Nhà nước (NHNN).
  \item \textbf{VIC} (bất động sản, ô tô): chi phí tài chính lớn và có nợ vay ngoại tệ, chịu cả hai kênh lãi suất
        và tỷ giá.
\end{itemize}

Chỉ số VN30 được đại diện bằng giá chứng chỉ quỹ ETF mô phỏng VN30 (chuỗi \texttt{VN30ETF}). Ngoài bốn cổ phiếu, dữ liệu gồm giá vàng thế giới XAU/USD, tỷ giá USD/VND và chỉ số đồng
đô la DXY. Hàm \texttt{market\_frame} lấy lịch giao dịch của HOSE làm trục thời gian và điền tiến các chuỗi nước
ngoài vào ngày giao dịch Việt Nam, nên mọi tài sản cùng có \SL{market/n_days} phiên chung từ \SL{market/range/0} tới
\SL{market/range/1}.

\bang{tab_market_desc}{Thống kê mô tả các tài sản trong giai đoạn nghiên cứu}{marketdesc}

\hinhpdf{mk_prices}{1.0}{Diễn biến giá điều chỉnh của bốn cổ phiếu, ETF VN30 và vàng quy đổi, chuẩn hóa về 100
ở đầu kỳ (thang log)}{mk_prices}

Bảng~\ref{tab:marketdesc} và Hình~\ref{fig:mk_risk} cho thấy lợi suất cao luôn đi cùng rủi ro cao. VIC có lợi suất
năm hóa cao nhất (\SL{market/describe/VIC/ann_retPct}\%) nhưng từng sụt \SL{market/describe/VIC/max_ddPct}\% từ
đỉnh; HPG sụt sâu nhất, \SL{market/describe/HPG/max_ddPct}\%, đúng tính chu kỳ của ngành thép. Vàng quy đổi ra
đồng có lợi suất \SL{market/describe/GOLD_VND/ann_retPct}\% mỗi năm, ngang VCB, nhưng mức sụt tối đa chỉ
\SL{market/describe/GOLD_VND/max_ddPct}\%, vì vậy vàng đóng vai trò tài sản phòng thủ trong danh mục của người Việt.
Tỷ giá USD/VND biến động chỉ \SL{market/describe/USDVND/ann_volPct}\% mỗi năm do NHNN điều hành trong biên độ, nhưng
chính vì ít biến động nên mỗi lần tỷ giá tăng nhanh đều là tin lớn với thị trường.

\hinhpdf{mk_risk}{0.92}{Lợi suất -- rủi ro năm hóa; diện tích bong bóng tỷ lệ với mức sụt giảm tối đa}{mk_risk}

\subsection{Tương quan giữa các tài sản}
Ma trận tương quan (Hình~\ref{fig:mk_corr}) được tính trên hai thang: lợi suất ngày và lợi suất cộng dồn 20 phiên.
Thang 20 phiên ít nhiễu hơn và gần với tầm nhìn của nhà đầu tư cá nhân. Ba nhóm quan hệ hiện ra rõ:

\begin{enumerate}
  \item Bốn cổ phiếu cùng tương quan dương với ETF VN30 (từ \SL{market/corr_20d/VIC/VN30ETF} với VIC tới
        \SL{market/corr_20d/VCB/VN30ETF} với VCB): rủi ro chung của thị trường chiếm phần lớn biến động.
  \item USD/VND tương quan \emph{âm} với cả bốn cổ phiếu (FPT \SL{market/corr_20d/USDVND/FPT}, HPG
        \SL{market/corr_20d/USDVND/HPG}): tỷ giá tăng thường đi cùng khối ngoại bán ròng và tâm lý rủi ro.
  \item Vàng thế giới gần như quyết định vàng quy đổi (\SL{market/corr_20d/XAU/GOLD_VND}), và ngược chiều với DXY
        (\SL{market/corr_20d/DXY/XAU}) vì vàng được định giá bằng USD.
\end{enumerate}

\hinhpdf{mk_corr}{1.0}{Ma trận tương quan Pearson của lợi suất ngày (trái) và lợi suất 20 phiên (phải)}{mk_corr}

Tương quan không đứng yên. Hình~\ref{fig:mk_rolling} vẽ tương quan trượt 252 phiên (khoảng một năm) của các cặp
tương ứng với cạnh trong KG: cặp DXY -- vàng luôn âm, còn cặp tỷ giá -- cổ phiếu dao động quanh 0 và đổi dấu theo
giai đoạn. Đây là cảnh báo đầu tiên rằng một cạnh có dấu cố định trong KG chỉ mô tả \emph{xu hướng cơ bản}, không
phải quy luật đúng ở mọi thời điểm.

\hinhpdf{mk_rolling}{1.0}{Tương quan trượt một năm của các cặp tương ứng với cạnh trong \texttt{VietnamFinanceKG}}{mk_rolling}

\subsection{Thị trường vàng Việt Nam: giá quy đổi và phần bù SJC}\label{sec:vang}
Giá vàng trong nước niêm yết theo \emph{lượng} (37,5 g), còn giá thế giới theo \emph{ounce troy} (31,1035 g). Giá
vàng thế giới quy đổi ra đồng, chưa gồm thuế phí, là

\begin{equation}\label{eq:vangquydoi}
  G_t \;=\; \frac{\mathrm{XAU}_t \times (\mathrm{USD/VND})_t \times 1{,}2057}{10^6}
  \quad\text{(triệu đồng/lượng)},
\end{equation}

\noindent với hệ số $1{,}2057 \approx 37{,}5 / 31{,}1035$ đổi giá mỗi ounce sang giá mỗi lượng. Giá vàng miếng SJC
bằng giá quy đổi cộng một \emph{phần bù} do cung cầu trong nước quyết định:
$P^{\mathrm{SJC}}_t = G_t + \Delta_t$. Ngày \SL{market/gold/last/date}, XAU/USD $=$ \SL{market/gold/last/xau}
USD/ounce và USD/VND $=$ \SL{market/gold/last/usdvnd}, nên $G_t \approx$ \SL{market/gold/last/gold_vnd} triệu
đồng/lượng. Công thức~\eqref{eq:vangquydoi} chính là công cụ \texttt{quy\_doi\_vang} của tác tử, và là lý do hai cạnh
\textsf{XAU\_USD} $\to$ \textsf{Vang\_Quy\_Doi}, \textsf{USD\_VND} $\to$ \textsf{Vang\_Quy\_Doi} trong KG đều mang
dấu dương.

\hinhtikz{t5_gold}{Cấu thành giá vàng miếng SJC: giá thế giới, tỷ giá và hệ số đổi đơn vị tạo nên giá quy đổi; phần
bù trong nước phụ thuộc chính sách}{t5_gold}

\hinhpdf{mk_gold}{1.0}{XAU/USD, USD/VND và giá vàng quy đổi theo~\eqref{eq:vangquydoi}}{mk_gold}

Phần bù $\Delta_t$ là đại lượng chịu tác động của chính sách, và là nơi KG cần tri thức pháp luật:

\begin{itemize}
  \item \textbf{Nghị định 24/2012/NĐ-CP}~\cite{nd24_2012} quy định Nhà nước độc quyền sản xuất vàng miếng, độc quyền
        xuất nhập khẩu vàng nguyên liệu để sản xuất vàng miếng; NHNN chọn SJC làm thương hiệu vàng miếng quốc gia.
        Từ năm 2014 NHNN gần như không nhập vàng nguyên liệu, nguồn cung bị giới hạn trong khi nhu cầu tích trữ tăng,
        nên phần bù có lúc lên tới 15 -- 18 triệu đồng/lượng đầu năm 2024.
  \item \textbf{Can thiệp năm 2024}: NHNN đấu thầu vàng miếng từ tháng 4, rồi bán trực tiếp qua bốn ngân hàng thương
        mại nhà nước và Công ty SJC từ tháng 6; phần bù thu hẹp còn khoảng 3 -- 5 triệu đồng/lượng.
  \item \textbf{Nghị định 232/2025/NĐ-CP}~\cite{nd232_2025} sửa đổi Nghị định 24, có hiệu lực từ tháng 10/2025,
        bãi bỏ cơ chế độc quyền sản xuất vàng miếng và cho phép doanh nghiệp, ngân hàng đủ điều kiện sản xuất vàng
        miếng, nhập khẩu vàng nguyên liệu theo hạn mức do NHNN quản lý.
\end{itemize}

Trong KG, ba sự kiện này là ba cạnh vào nút \textsf{Chenh\_Lech\_SJC}: Nghị định 24 làm phần bù tăng (dấu $+$),
hai biện pháp còn lại làm phần bù giảm (dấu $-$). Vàng nhẫn 9999 không chịu cơ chế độc quyền nên chỉ nối với giá
quy đổi, không nối với phần bù.

% ---------------------------------------------------------------------
\section{Đồ thị tri thức tài chính \texttt{VietnamFinanceKG}}\label{sec:kgtc}
\subsection{Bản thể học}
Bản thể học tài chính nhỏ hơn bản thể học y tế nhưng có cấu trúc khác: \SL{kg/finance/classes} lớp và chỉ
\SL{kg/finance/predicates} vị từ. Lớp gốc \textsf{FinancialEntity} tách thành yếu tố vĩ mô (\textsf{MacroFactor},
với ba lớp con \textsf{PolicyTool}, \textsf{GlobalFactor}, \textsf{MarketFlow}), tài sản (\textsf{Asset} $\to$
\textsf{Commodity} $\to$ \textsf{GoldAsset}, \textsf{Stock}, \textsf{Index}), văn bản pháp luật, sự kiện, chỉ báo,
ngành và sàn giao dịch. Hai vị từ quan trọng nhất mang \emph{dấu}: \textsf{TAC\_DONG\_CUNG\_CHIEU} ($+1$, biến
nguồn tăng thì biến đích có xu hướng tăng) và \textsf{TAC\_DONG\_NGUOC\_CHIEU} ($-1$). Bốn vị từ còn lại là quan hệ
cấu trúc (thành viên chỉ số, ngành, sửa đổi văn bản, phạm vi điều chỉnh) và không tham gia lan truyền tác động.
\textsf{THUOC\_VE} được khai báo bắc cầu, nên suy diễn RDFS (Hình~\ref{fig:suydien}) tự sinh ``FPT thuộc về HOSE'' từ
hai cạnh ``FPT thuộc VN30'' và ``VN30 thuộc HOSE''.

\maNguon{kg_finance__finance_ontology}{Bản thể học tài chính: \SL{kg/finance/classes} lớp và
\SL{kg/finance/predicates} vị từ, hai vị từ nhân quả mang dấu}

\subsection{Thực thể và quan hệ nhân quả}
KG có \SL{kg/finance/nodes} thực thể và \SL{kg/finance/edges} cạnh, trong đó
\SL{kg/finance/relations/TAC_DONG_CUNG_CHIEU} cạnh cùng chiều và \SL{kg/finance/relations/TAC_DONG_NGUOC_CHIEU}
cạnh ngược chiều. Cả \SL{kg/finance/with_source} cạnh đều có thuộc tính \texttt{source} trỏ về tài liệu trong kho
văn bản (\texttt{f01} -- \texttt{f09}) và thuộc tính \texttt{mechanism} giải thích cơ chế bằng một câu, nên mỗi kết
luận của tác tử đều truy được về văn bản gốc. Thực thể có chuỗi giá thật mang thuộc tính \texttt{series}, nối KG với
dữ liệu thị trường; nhờ đó Mục~\ref{sec:kiemchung} kiểm chứng được dấu cạnh. Hình~\ref{fig:t5_causal} tách đồ thị
nhân quả thành hai kênh để dễ đọc, Hình~\ref{fig:kg_finance_graph} vẽ toàn bộ KG.

\hinhtikz{t5_causal}{Đồ thị nhân quả của \texttt{VietnamFinanceKG}: (a) kênh vàng, (b) kênh tỷ giá -- lãi suất --
cổ phiếu; xanh liền là cùng chiều, đỏ đứt là ngược chiều}{t5_causal}

\maDai{kg_finance__build_finance_kg}{Xây \texttt{VietnamFinanceKG}: thực thể kèm bí danh, chuỗi giá, và các cạnh
nhân quả có nguồn, có cơ chế}

\hinhpdf{kg_finance_graph}{1.0}{Toàn cảnh \texttt{VietnamFinanceKG}: \SL{kg/finance/nodes} thực thể,
\SL{kg/finance/edges} quan hệ; màu nút theo lớp, màu cạnh theo dấu tác động}{kg_finance_graph}

Khi xuất ra RDF, KG có \SL{kg/finance/rdf_triples} bộ ba; bộ suy diễn RDFS thêm \SL{kg/finance/inferred} bộ ba
(kiểu kế thừa từ lớp cha, quan hệ bắc cầu) để tổng số đạt \SL{kg/finance/rdf_after}. Chẳng hạn ``FPT là một
\textsf{Stock}'' kéo theo ``FPT là một \textsf{Asset}'' và ``FPT là một \textsf{FinancialEntity}''. Nhờ vậy truy vấn SPARQL hỏi ``mọi tài sản chịu tác động ngược chiều của tỷ giá'' không
cần liệt kê từng lớp con.

% ---------------------------------------------------------------------
\section{Suy luận tác động đa bước}\label{sec:multihop}
\subsection{Dấu của đường đi và tác động ròng}
Với một biến nguồn $u$ (chẳng hạn lãi suất Fed) và một biến đích $v$ (vàng SJC), hàm \texttt{impact\_paths} liệt kê
mọi đường đơn $\pi = (u = x_0 \to x_1 \to \dots \to x_k = v)$ dài không quá 4 bước trên đồ thị con nhân quả. Dấu của
đường là tích dấu các cạnh, và tác động ròng là tổng dấu các đường có trọng số nghịch đảo độ dài:

\begin{equation}\label{eq:rong}
  s(\pi) \;=\; \prod_{j=1}^{k} s(x_{j-1}, x_j), \qquad
  N(u, v) \;=\; \sum_{\pi:\, u \rightsquigarrow v} \frac{s(\pi)}{|\pi|}.
\end{equation}

Trọng số $1/|\pi|$ thể hiện giả định rằng tác động suy giảm qua mỗi bước trung gian: một kênh trực tiếp đáng tin hơn
một chuỗi bốn mắt xích. Khi mọi đường cùng dấu, kết luận là ``tích cực'' hoặc ``tiêu cực''; khi các đường mâu thuẫn,
bộ suy luận \emph{không} gộp thành một chiều mà trả về ``hỗn hợp'' kèm số kênh mỗi dấu, để tác tử nói rõ kết quả phụ
thuộc kênh truyền dẫn nào mạnh hơn.

\maNguon{kg_finance__impact_paths}{Suy luận đa bước: liệt kê đường nhân quả, tính dấu từng đường và tác động
ròng~\eqref{eq:rong}}

\hinhtikz{t5_multihop}{Ba đường nhân quả từ lãi suất Fed tới vàng SJC mà \texttt{impact\_paths} tìm được, dấu của
từng đường và tác động ròng}{t5_multihop}

\bang{tab_multihop}{Kết quả suy luận đa bước cho bảy cặp nguồn -- đích}{multihop}

Bảng~\ref{tab:multihop} có hai trường hợp ``hỗn hợp'' đáng phân tích:

\begin{itemize}
  \item \textbf{Fed $\to$ vàng SJC}: hai đường âm (Fed $\to$ XAU $\to$ quy đổi $\to$ SJC, và Fed $\to$ DXY $\to$ XAU
        $\to$ quy đổi $\to$ SJC) và một đường dương (Fed $\to$ DXY $\to$ USD/VND $\to$ quy đổi $\to$ SJC). Đường dương
        có thật về kinh tế: Fed tăng lãi suất làm USD mạnh, tỷ giá tăng, và giá vàng tính bằng đồng được đỡ một phần.
        Tác động ròng $N = -1/3 - 1/4 + 1/4 =$ \SL{kg/demo_finance/Lai_Suat_Fed->Vang_SJC/net}: chiều chi phí cơ hội
        của vàng thắng, nhưng nhà đầu tư trong nước chịu giảm ít hơn nhà đầu tư nắm vàng tính bằng USD.
  \item \textbf{Lãi suất NHNN $\to$ HPG}: một đường trực tiếp âm (lãi suất cao làm giảm cầu xây dựng) và một đường
        hai bước dương (lãi suất VND cao làm giảm áp lực tỷ giá, mà tỷ giá thấp có lợi cho doanh nghiệp nhập nguyên
        liệu). Ròng $-1 + 1/2 =$ \SL{kg/demo_finance/Lai_Suat_NHNN->HPG/net}.
\end{itemize}

\subsection{Lan truyền cú sốc}
Câu hỏi ngược lại, ``một cú sốc ở $u$ ảnh hưởng tới những gì'', được trả lời bằng hàm \texttt{downstream}: duyệt
theo chiều rộng từ $u$ tối đa 3 bước, mỗi thực thể nhận dấu của đường ngắn nhất tới nó. Hình~\ref{fig:kg_shock} và
Bảng~\ref{tab:shock} là kết quả cho cú sốc DXY tăng: mười thực thể bị ảnh hưởng. FPT là cổ phiếu duy nhất có lợi
(qua tỷ giá); VN30, HPG, VIC và cả bốn nút vàng đều chịu chiều giảm; VCB không nằm trong danh sách vì KG không có
kênh nối VCB với tỷ giá.

\maNguon{kg_finance__downstream}{Lan truyền cú sốc theo chiều rộng, giữ dấu của đường ngắn nhất}

\hinhpdf{kg_shock}{0.92}{Lan truyền cú sốc ``DXY tăng'' trong \texttt{VietnamFinanceKG}: xanh là tăng, đỏ là
giảm}{kg_shock}

\bang{tab_shock}{Chiều tác động lên các thực thể khi DXY tăng, theo đường ngắn nhất}{shock}

\subsection{Kiểm chứng dấu cạnh bằng dữ liệu}\label{sec:kiemchung}
Một KG tài chính viết tay dễ chứa những ``hiểu biết thông thường'' sai. Bài kiểm chứng mọi cạnh nhân quả có chuỗi
giá ở cả hai đầu: tính tương quan Pearson giữa lợi suất 20 phiên của biến nguồn và biến đích trên toàn giai đoạn, rồi
so dấu với dấu lý thuyết. Kết quả (Bảng~\ref{tab:edgeev}, Hình~\ref{fig:mk_edges}): \SL{dem/edges_agree} trên
\SL{dem/edges_checked} cạnh khớp dấu. Cạnh \textsf{XAU\_USD} $\to$ \textsf{Vang\_Quy\_Doi} gần như tất định
(\SL{market/corr_20d/XAU/GOLD_VND}); cạnh DXY $\to$ XAU mạnh và đúng chiều; các cạnh tỷ giá $\to$ HPG, VIC yếu
nhưng đúng dấu.

\bang{tab_edges_evidence}{Đối chiếu dấu lý thuyết của cạnh KG với tương quan thực nghiệm}{edgeev}

\hinhpdf{mk_edges}{1.0}{Tương quan thực nghiệm (thanh) và dấu lý thuyết trong KG (hình thoi) của bảy cạnh có chuỗi
giá ở hai đầu}{mk_edges}

Cạnh không khớp là USD/VND $\to$ FPT: lý thuyết cho dấu dương, dữ liệu cho tương quan cùng kỳ
\SL{market/corr_20d/USDVND/FPT}. Sai lệch này không có nghĩa cạnh sai, mà cho thấy hai cơ chế chạy ở hai thang thời
gian khác nhau. Trong ngắn hạn, tỷ giá tăng nhanh đi cùng dòng vốn ngoại rút khỏi thị trường và tâm lý né rủi ro, kéo
mọi cổ phiếu lớn giảm cùng lúc, kể cả FPT. Lợi ích cơ bản của doanh thu ngoại tệ chỉ hiện ra khi doanh nghiệp công
bố kết quả kinh doanh, tức là vài tháng sau. Tương quan trễ (Hình~\ref{fig:mk_lag}) xác nhận cách giải thích này:
tương quan giữa lợi suất 20 phiên của tỷ giá hôm nay và của FPT $k$ phiên sau tăng dần theo $k$, đổi dấu quanh
$k = 20$ và đạt cực đại \SL{dem/lag_corr_FPT} ở độ trễ \SL{dem/lag_peak_FPT} phiên, trong khi với HPG, đường tương
quan chỉ tiến về 0 mà không đổi dấu đáng kể.

\hinhpdf{mk_lag}{1.0}{Tương quan giữa lợi suất 20 phiên của USD/VND và của cổ phiếu sau $k$ phiên}{mk_lag}

\nhanxet{Cạnh KG mã hóa cơ chế cơ bản, còn tương quan cùng kỳ đo phản ứng của dòng tiền. Bài giữ nguyên cạnh
USD/VND $\to$ FPT vì có căn cứ văn bản; đây là ví dụ cho thấy KG tài chính cần được đối chiếu với dữ liệu thay vì chỉ
dựa vào hiểu biết thông thường, và là lý do Chương~\ref{chap:ketluan} đề xuất gắn tương quan thực nghiệm cùng độ trễ
vào thuộc tính của từng cạnh. Một lưu ý về phương pháp: lợi suất 20 phiên tính trên cửa sổ chồng lấn nên các quan sát liền nhau phụ
thuộc nhau; độ lớn tương quan dùng được để so chiều, nhưng không nên đọc như kiểm định ý nghĩa thống kê.}

% ---------------------------------------------------------------------
\section{Dự báo chuỗi thời gian bằng BiLSTM}\label{sec:bilstm}
\subsection{Phát biểu bài toán}
Mô hình không dự báo trực tiếp giá mà dự báo \emph{lợi suất log} phiên kế tiếp, $r_{t+1} = \ln(P_{t+1}/P_t)$,
rồi suy ra giá $\hat P_{t+1} = P_t\, e^{\hat r_{t+1}}$. Chuỗi giá không dừng: một mô hình hồi quy trên giá đạt $R^2$
rất cao chỉ bằng cách chép giá hôm qua, và con số đó không nói gì về khả năng dự báo. Lợi suất gần dừng hơn, và cho
phép so sánh với mốc khó nhất: \emph{dự báo ngây thơ} ``giá không đổi'', tức $\hat r_{t+1} = 0$. Với mốc này,
$\mathrm{RMSE}_{\text{ngây thơ}} = \sqrt{\tfrac{1}{n}\sum r_{t+1}^2}$, và chỉ số chính của bài là tỉ số
$\mathrm{RMSE}/\mathrm{RMSE}_{\text{ngây thơ}}$: dưới 1 mới là có năng lực dự báo.

Mỗi phiên có 7 đặc trưng với cổ phiếu (6 với vàng): lợi suất log của chính tài sản; lợi suất của bốn biến ngoại sinh
USD/VND, XAU, DXY, ETF VN30 (bỏ XAU khi dự báo vàng quy đổi vì XAU là thành phần của chính nó); độ lệch của giá so
với trung bình trượt 20 phiên; và độ biến động 20 phiên. Mỗi mẫu là một cửa sổ
\SL{forecast/config/window} phiên liên tiếp, tức một tenxơ $30 \times 7$.

\maNguon{data__forecast_features}{Đặc trưng dự báo: lợi suất của đích và biến ngoại sinh, độ lệch MA20, độ biến
động 20 phiên; \texttt{keep\_last} giữ phiên mới nhất khi suy luận}

\maNguon{data__make_windows}{Cắt chuỗi đặc trưng thành các cửa sổ trượt cho mô hình tuần tự}

Dữ liệu được chia \emph{theo thời gian}, không xáo trộn (Hình~\ref{fig:t5_split}): huấn luyện trước
\SL{forecast/config/val_start}, kiểm định năm 2024, kiểm thử từ \SL{forecast/config/test_start} tới hết dữ liệu.
Trung bình và độ lệch chuẩn để chuẩn hóa chỉ tính trên tập huấn luyện. Với FPT, ba tập có \SL{forecast/FPT/n/train},
\SL{forecast/FPT/n/val} và \SL{forecast/FPT/n/test} cửa sổ.

\hinhtikz{t5_split}{Chia dữ liệu theo thời gian và cửa sổ trượt 30 phiên: mô hình không bao giờ nhìn thấy tương lai
của mẫu nó dự báo}{t5_split}

\subsection{Kiến trúc LSTM hai chiều}
LSTM~\cite{hochreiter1997lstm} giải quyết bài toán gradient biến mất của RNN bằng một ô nhớ $c_t$ và ba cổng điều
khiển luồng thông tin. Với đầu vào $x_t$ và trạng thái ẩn trước $h_{t-1}$:

\begin{align}
  f_t &= \sigma\!\left(W_f [h_{t-1}, x_t] + b_f\right) & &\text{cổng quên} \label{eq:lstm-f}\\
  i_t &= \sigma\!\left(W_i [h_{t-1}, x_t] + b_i\right), \quad
  \tilde c_t = \tanh\!\left(W_c [h_{t-1}, x_t] + b_c\right) & &\text{cổng vào, ứng viên} \\
  c_t &= f_t \odot c_{t-1} + i_t \odot \tilde c_t & &\text{cập nhật ô nhớ} \\
  o_t &= \sigma\!\left(W_o [h_{t-1}, x_t] + b_o\right), \quad h_t = o_t \odot \tanh(c_t) & &\text{cổng ra} \label{eq:lstm-o}
\end{align}

LSTM hai chiều~\cite{schuster1997brnn} chạy thêm một LSTM ngược thời gian trên \emph{cùng cửa sổ} và ghép hai trạng
thái, $h_t = [\overrightarrow{h}_t ; \overleftarrow{h}_t]$. Vì cửa sổ chỉ gồm các phiên đã qua, chiều ngược đọc từ
phiên $t$ lùi về phiên $t-29$ và không làm rò rỉ tương lai; tác dụng của nó là để trạng thái cuối mang thông tin của
cả đầu lẫn cuối cửa sổ. Mạng gồm hai tầng BiLSTM 64 chiều mỗi hướng, lấy trạng thái ở bước cuối (128 chiều) qua
Dropout và hai lớp tuyến tính $128 \to 32 \to 1$; tổng cộng 140.865 tham số với 7 đặc trưng, so với 54.081 tham số
của LSTM một chiều cùng cấu hình.

\hinhtikz{t5_bilstm3d}{Bộ dự báo BiLSTM: tenxơ cửa sổ $30 \times 7$, hai tầng LSTM xuôi và ngược thời gian, đầu
tuyến tính cho lợi suất phiên kế tiếp và MC Dropout cho khoảng dự báo}{t5_bilstm3d}

\maNguon{dl_models__SeqForecaster}{Mô hình tuần tự dùng chung cho BiLSTM và LSTM một chiều}

\bang{tab_fc_config}{Cấu hình bộ dự báo}{fcconfig}

\subsection{Huấn luyện}
Hàm mất mát là Huber~\cite{huber1964robust} với $\delta = 1$ trên lợi suất đã chuẩn hóa:
\begin{equation}
  \mathcal{L}_\delta(e) = \begin{cases} \tfrac12 e^2 & |e| \le \delta \\ \delta\left(|e| - \tfrac12\delta\right) &
  |e| > \delta \end{cases}
\end{equation}
Lợi suất có đuôi dày: vài phiên giảm sàn hay tăng trần có thể chi phối MSE, trong khi Huber chỉ phạt tuyến tính các
sai số lớn. Tối ưu bằng AdamW~\cite{loshchilov2019adamw} (tốc độ học $10^{-3}$, suy giảm trọng số $10^{-4}$), lô 128,
cắt chuẩn gradient ở 1,0 để tránh bùng nổ gradient trong LSTM, tối đa 80 epoch và dừng sớm khi mất mát kiểm định
không giảm sau 10 epoch. Huấn luyện chạy trên GPU, mỗi mô hình mất vài giây
(Bảng~\ref{tab:forecast}). Đối chứng gồm LSTM một chiều cùng cấu hình, hồi quy Ridge trên 30 trễ của mọi đặc
trưng, và dự báo ngây thơ.

\maDai{train_forecast__train_rnn}{Vòng huấn luyện: Huber, AdamW, cắt gradient, dừng sớm theo mất mát kiểm định}

\maNguon{train_forecast__evaluate}{Đánh giá trên tập kiểm thử: RMSE lợi suất, tỉ số với dự báo ngây thơ, MAE giá,
độ chính xác hướng}

\hinhpdf{fc_history}{1.0}{Đường cong học của BiLSTM cho FPT và vàng quy đổi}{fc_history}

Đường cong học (Hình~\ref{fig:fc_history}) cho thấy dấu hiệu đầu tiên của kết quả: mất mát huấn luyện giảm đều,
còn mất mát kiểm định gần như đi ngang hoặc tăng ngay từ những epoch đầu. Với HPG, VIC và vàng quy đổi, trọng số tốt
nhất là của epoch thứ nhất và quá trình dừng ở epoch 11; với FPT, mất mát kiểm định tốt nhất (epoch 22) chỉ thấp hơn
epoch đầu khoảng 2\%. Phần mô hình học thêm từ tập huấn luyện gần như không tổng quát hóa sang năm 2024.

\subsection{Kết quả và diễn giải}
\bang{tab_forecast}{Kết quả dự báo trên tập kiểm thử của năm tài sản và bốn mô hình}{forecast}

\haihinh{fc_ratio}{Tỉ số RMSE so với dự báo ngây thơ}{fc_diracc}{Tỉ lệ đoán đúng chiều tăng/giảm}{So sánh BiLSTM,
LSTM và Ridge trên năm tài sản}{fc_compare}

Kết quả cần được nói thẳng. Chỉ \SL{dem/fc_beat} trên \SL{dem/fc_models} mô hình có tỉ số RMSE dưới 1, và mô hình
tốt nhất (LSTM trên vàng quy đổi) chỉ giảm sai số \SL{dem/fc_best_gainPct}\% so với dự báo ngây thơ, một mức không có
ý nghĩa thực tế. BiLSTM của FPT, mô hình phức tạp nhất trên tài sản biến động mạnh nhất trong nhóm, kém mốc ngây thơ
tới tỉ số \SL{dem/fc_worst_ratio}. Ridge trên 30 trễ kém dự báo ngây thơ ở mọi tài sản: với hơn 200 đặc trưng tuyến
tính, mô hình khớp nhiễu của tập huấn luyện. Tỉ lệ đoán đúng chiều dao động quanh 50\%, cao nhất là LSTM trên vàng
(\SL{forecast/GOLD_VND/lstm/dir_accPct}\%) và thấp nhất là hai mô hình LSTM trên VIC (khoảng 41\%, tức tệ hơn tung
đồng xu). Dự báo ngây thơ không có cột đúng chiều
vì lợi suất 0 không có dấu.

\hinhpdf{fc_test}{1.0}{Giá thực tế và giá dự báo một phiên của BiLSTM trên tập kiểm thử}{fc_test}

Hình~\ref{fig:fc_test} dễ gây hiểu lầm nếu đọc thiếu cẩn thận: đường dự báo bám sát đường thực tế, nhưng đó là vì
$\hat P_{t+1} = P_t e^{\hat r_{t+1}}$ và lợi suất một phiên chỉ vài phần trăm, nên bất kỳ dự báo nào gần giá hôm qua
cũng trông khớp. Hình~\ref{fig:fc_scatter} đặt lợi suất dự báo cạnh lợi suất thực tế và cho thấy hai kiểu thất bại
khác nhau. Với vàng quy đổi, BiLSTM \emph{co hẳn về 0}: độ lệch chuẩn của dự báo chỉ \SL{dem/fc_std_pred_GOLD_VND}\% so với
\SL{dem/fc_std_true_GOLD_VND}\% của thực tế, tức mô hình tự học ra dự báo ngây thơ, và tỉ số RMSE vì vậy sát 1. Với
FPT, BiLSTM làm điều ngược lại: dự báo dao động mạnh (độ lệch chuẩn \SL{dem/fc_std_pred_FPT}\% so với
\SL{dem/fc_std_true_FPT}\%) và lệch dương trung bình \SL{dem/fc_mean_pred_FPT}\% mỗi phiên, nhưng các
dao động đó không trùng chiều với thực tế, nên mỗi lần đoán mạnh là một lần cộng thêm sai số; đó là lý do tỉ số RMSE
của nó tệ nhất bảng. Cả hai kiểu đều dẫn tới cùng kết luận: lợi suất ngày mai gần như không dự đoán được từ thông tin
công khai của 30 phiên trước, và mô hình nào ``dám'' dự báo xa khỏi 0 thì bị phạt.

\hinhpdf{fc_scatter}{1.0}{Lợi suất dự báo và lợi suất thực tế trên tập kiểm thử: FPT dự báo mạnh nhưng sai chiều,
vàng quy đổi co về 0}{fc_scatter}

Kết quả này phù hợp với giả thuyết thị trường hiệu quả dạng yếu~\cite{fama1970efficient}: nếu giá đã phản ánh mọi
thông tin trong lịch sử giá, không mô hình nào chỉ dùng lịch sử giá có thể dự báo lợi suất tốt hơn bước ngẫu nhiên.
Thị trường Việt Nam được xem là kém hiệu quả hơn thị trường phát triển, nhưng ở tầm một phiên và với các cổ phiếu
thanh khoản cao nhất VN30, phần dự báo được, nếu có, nhỏ hơn nhiễu. Bài không coi đây là thất bại cần che đi, mà là
một kết quả thực nghiệm quyết định cách tác tử dùng mô hình.

\subsection{Khoảng dự báo cho tác tử}
Công cụ \texttt{du\_bao\_bilstm} của tác tử không chỉ trả một con số. Lớp \texttt{Forecaster} lấy cửa sổ 30 phiên
mới nhất và chạy 200 lượt với Dropout bật, mỗi lượt là một mẫu từ phân phối hậu nghiệm xấp xỉ của trọng số theo MC
Dropout~\cite{gal2016dropout}. Tuy nhiên MC Dropout chỉ đo \emph{bất định của mô hình}; với chuỗi lợi suất, phần này
rất nhỏ so với nhiễu vốn có của giá. Phiên bản đầu của công cụ dùng riêng MC Dropout và báo ``xác suất tăng 100\%''
với khoảng 90\% rộng chưa tới $\pm 0{,}2\%$ quanh giá hiện tại, một kết luận tự tin sai hoàn toàn. Bản cuối cộng vào
mỗi mẫu một nhiễu phần dư có độ lệch chuẩn bằng RMSE trên tập kiểm thử:

\begin{equation}\label{eq:khoangdubao}
  r^{(s)} = f_{\theta^{(s)}}(\mathbf{x}_{t-29:t}) + \varepsilon^{(s)}, \quad
  \varepsilon^{(s)} \sim \mathcal{N}\!\left(0, \mathrm{RMSE}^2_{\text{kiểm thử}}\right), \quad s = 1, \dots, 200,
\end{equation}

\noindent rồi lấy phân vị 5\% và 95\% của $P_t e^{r^{(s)}}$ làm khoảng dự báo 90\% và tỉ lệ $r^{(s)} > 0$ làm xác suất
tăng. Bảng~\ref{tab:fclive} đặt hai nguồn bất định cạnh nhau: với FPT, độ lệch chuẩn của MC Dropout chỉ
\SL{forecast_live/FPT/mc_std_pct}\% trong khi RMSE kiểm thử là \SL{forecast_live/FPT/rmse_pct}\%, tức nhiễu của dữ
liệu lớn hơn bất định của mô hình hơn hai mươi lần. Xác suất tăng sau khi sửa nằm trong khoảng 50 -- 60\%, đúng với
mức hiểu biết thật của mô hình.

\maNguon{dl_models__Forecaster.predict}{Công cụ dự báo của tác tử: MC Dropout cộng nhiễu phần dư cho khoảng dự báo
90\% và xác suất tăng}

\bang{tab_fc_live}{Dự báo phiên tới của năm tài sản: bất định mô hình so với sai số dự báo thực tế}{fclive}

Văn bản công cụ trả cho tác tử luôn đi kèm hai chỉ số kiểm thử: tỉ lệ đoán đúng chiều và tỉ số RMSE so với dự báo
ngây thơ, cùng câu ``không vượt dự báo giữ nguyên giá'' khi tỉ số không dưới 1. Như vậy câu trả lời cuối cùng của tác
tử tự mang theo mức độ tin cậy của chính nó.

% ---------------------------------------------------------------------
\section{Tình huống minh họa: tác tử trả lời nhà đầu tư}\label{sec:tc-tinhhuong}
\subsection{Câu hỏi tác động vĩ mô}
Câu hỏi ``Tỷ giá USD/VND tăng mạnh thì cổ phiếu FPT và giá vàng SJC ảnh hưởng thế nào?'' cần cả ba lớp. Liên kết
thực thể nhận ra \textsf{USD\_VND}, \textsf{FPT} và \textsf{Vang\_SJC}; bộ lập kế hoạch gọi \texttt{kg\_tai\_chinh}
để suy luận hai cặp tác động, \texttt{du\_lieu\_thi\_truong} để lấy giá FPT gần nhất, \texttt{quy\_doi\_vang} để tính
giá quy đổi theo~\eqref{eq:vangquydoi}, và \texttt{truy\_xuat\_rag} để lấy đoạn văn bản giải thích cơ chế. Khung
dưới đây là vết suy luận thật của tác tử, do \texttt{build.py} sinh từ \texttt{agent.json}.

\begin{vetsuyluan}{Vết suy luận ReAct thật cho câu hỏi tài chính (kịch bản 9)}
\input{generated/trace_finance}
\end{vetsuyluan}

\begin{vetsuyluan}{Câu trả lời cuối của tác tử (bản tổng hợp theo khuôn, có trích dẫn [n])}
\input{generated/answer_finance}
\end{vetsuyluan}

Bản nháp của Qwen2.5-1.5B cho câu hỏi này không sai số liệu nào (trung thực số học 1,00) nhưng không ghi số hiệu
bằng chứng [n] nào, nên bộ kiểm chứng từ chối và trả bản tổng hợp theo khuôn; Mục~\ref{sec:agent-ketqua} phân tích
hiện tượng này trên toàn bộ kịch bản. Bản cuối đặt cạnh nhau hai loại bằng chứng: kết luận suy luận từ KG (FPT và
vàng SJC cùng chịu tác động cùng chiều của tỷ giá, mỗi kết luận kèm kênh truyền dẫn và cơ chế) và số liệu tính từ dữ
liệu (giá FPT, giá vàng quy đổi), mỗi dòng mang số hiệu [n] trỏ về bằng chứng gốc. Hình~\ref{fig:ui_finance} là cùng
câu hỏi trên Web App.

\anhUI{ui_finance}{Trợ lý tài chính trên Web App: câu hỏi tác động tỷ giá, vết suy luận và câu trả lời có trích
dẫn}{ui_finance}

\subsection{Câu hỏi dự báo}
Với ``Dự báo giá cổ phiếu FPT phiên tới?'', bộ lập kế hoạch nhận ra ý định dự báo và gọi \texttt{du\_bao\_bilstm} sau
\texttt{du\_lieu\_thi\_truong} (Hình~\ref{fig:ui_finance_forecast}). Câu trả lời nêu giá dự báo, khoảng dự báo 90\%,
xác suất tăng, và ngay sau đó là hai chỉ số kiểm thử cho thấy mô hình không vượt dự báo ngây thơ. Một nhà đầu tư đọc
câu trả lời này biết chính xác nên tin con số tới đâu, điều mà một câu trả lời chỉ có ``giá dự kiến 63.315 đồng''
không làm được.

\anhUI{ui_finance_forecast}{Câu hỏi dự báo FPT: tác tử gọi BiLSTM và trả kèm khoảng dự báo, xác suất tăng và độ tin
cậy trên tập kiểm thử}{ui_finance_forecast}

\subsection{Khám phá đồ thị tri thức trên giao diện}
Trang ``Đồ thị tri thức'' của Web App vẽ \texttt{VietnamFinanceKG} bằng SVG thuần, tô màu nút theo lớp, cạnh ngược
chiều màu đỏ (Hình~\ref{fig:ui_kg_finance}). Người dùng nhấn vào một nút để xem lớp của nó và mọi cạnh vào ra kèm
tài liệu nguồn; đây là cách nhanh nhất để kiểm tra vì sao tác tử kết luận một chiều tác động.

\anhUI{ui_kg_finance}{Trang khám phá \texttt{VietnamFinanceKG} trên Web App}{ui_kg_finance}

% ---------------------------------------------------------------------
\section{Tổng kết chương}
\begin{itemize}
  \item KG tài chính với \SL{kg/finance/nodes} thực thể và \SL{kg/finance/edges} cạnh có dấu, có nguồn, đủ để suy
        luận đa bước và lan truyền cú sốc; khi các kênh mâu thuẫn, bộ suy luận trả ``hỗn hợp'' thay vì ép một chiều.
  \item \SL{dem/edges_agree} trên \SL{dem/edges_checked} cạnh kiểm chứng được khớp dấu với dữ liệu; cạnh tỷ giá
        $\to$ FPT chỉ khớp khi xét độ trễ khoảng \SL{dem/lag_peak_FPT} phiên.
  \item BiLSTM không thắng có ý nghĩa dự báo ngây thơ ở tầm một phiên, phù hợp giả thuyết thị trường hiệu quả dạng
        yếu; giá trị của nó với tác tử nằm ở khoảng dự báo trung thực và các chỉ số tin cậy đi kèm, không ở con số
        điểm.
  \item Khoảng dự báo chỉ dựa trên MC Dropout tự tin sai; cộng nhiễu phần dư theo~\eqref{eq:khoangdubao} đưa xác
        suất tăng về đúng mức 50 -- 60\%.
\end{itemize}
